%% notes-tex/multimodal/043-training-design/sections/4-objective.tex
%% SOURCE: torchcell/losses/hl_gauss.py (HLGaussLoss, fit_standardization),
%% torchcell/losses/distributional.py (pinball), torchcell/models/cell_model_043.py
%% (ReadoutLoss043), train_043.py (SPECS, transform_target, loss_sums, combine_loss,
%% balanced_batches, batch_loader, predict, evaluate, OptimConfig), conf/loss/*.yaml,
%% conf/weighting/*.yaml; batch shares through tables/t2-batch-share-rows.tex and
%% tables/macros.tex (training_design_043_tables.py); standardization sha256 from the job 6136 log.

\section{Objective, step by step\secstatus{todo}}
\label{sec:objective}

\subsection{Targets\secstatus{todo}}
\label{sec:targets}

\notesec{Step 1, the transform}
Each dataset's released value is first put on a scale where differences are comparable,
$\tilde y_{b,i} = \operatorname{tr}_s(y_{b,i})$: as released for Kemmeren (already a
$\log_2$ ratio to wild type) and Vaishnav (a sort-seq level); the $\log_2$ ratio to the
reference for Messner; the difference to the reference for Zelezniak; $\log_2(y+1)$ of TPM
for Caudal; $\log_2 y$ of positive values for Muenzner protein and turnover; and
$\log_2(10^{6} y / \text{UMI total} + 1)$ for the Boocock and N'Guessan counts, where a
gene a record omits is a released zero. The primer's table of output scales gives the
reasons per dataset.

\notesec{Step 2, standardization per dataset}
On the training rows of dataset $s$ only,
\begin{equation}
z_{b,i} = \frac{\tilde y_{b,i} - \eta_{s,i} - \mu_s}{\sigma_s},
\qquad
m_{b,i} = \mathbf{1}\big[z_{b,i} \text{ is finite}\big],
\label{eq:standardize}
\end{equation}
where $\eta_{s,i}$ is the mean of gene $i$ over the training rows that measure it for the
five panel datasets centered per gene (Caudal, Muenzner protein and turnover, Boocock,
N'Guessan) and zero for the rest, and $\mu_s$, $\sigma_s$ are the mean and standard
deviation pooled over every centered training value of $s$. A gene with no training value
has no center, so its $z$ is missing rather than zero. The fitted statistics are one
pydantic record pinned by sha256 (\texttt{0d5df0f0\ldots} in the overnight runs, job 6136
log). Counts never detected in a Boocock or N'Guessan training row are not targets.

\subsection{Elementwise losses\secstatus{todo}}
\label{sec:elementwise}

\notesec{Step 3, HL-Gauss bins}
The main loss (decision O4 of the primer) bins $z$ into $n_{\mathrm{bin}} = 32$ bins on
$[-4, 4]$ of width $\Delta = 0.25$, with interior edges $\upsilon_u = -4 + u\Delta$ for $u =
1, \ldots, 31$, the outer edges opened to $\upsilon_0 = -\infty$ and $\upsilon_{32} =
+\infty$, and nominal centers $\bar\upsilon_u = -4 + (u - \tfrac12)\Delta$.

\notesec{Step 4, the Gaussian histogram target}
A measured $z$ becomes the mass a normal $\mathcal{N}(z, \sigma_{\mathrm{HL}}^2)$ puts in
each bin, with $\sigma_{\mathrm{HL}} = 0.75\Delta = 0.1875$ and $\Phi$ the standard normal
distribution function,
\begin{equation}
\omega_u(z) = \Phi\!\Big(\frac{\upsilon_u - z}{\sigma_{\mathrm{HL}}}\Big) - \Phi\!\Big(\frac{\upsilon_{u-1} - z}{\sigma_{\mathrm{HL}}}\Big),
\qquad u = 1, \ldots, n_{\mathrm{bin}},
\qquad \sum_{u} \omega_u(z) = 1.
\label{eq:hltarget}
\end{equation}
Because the edge bins are open, a target beyond $\pm 4$ still has a histogram that sums to
one; its mass sits in the edge bin.

\notesec{Step 5, the cross-entropy}
The head's $n_{\mathrm{out}} = 32$ outputs are logits,
\begin{equation}
\hat\omega_u = \frac{\exp(\hat o_u)}{\sum_{u'} \exp(\hat o_{u'})},
\qquad
\mathcal{L}^{\mathrm{HL}}(\hat{\mathbf{o}}, z) = -\sum_{u=1}^{n_{\mathrm{bin}}} \omega_u(z) \log \hat\omega_u.
\label{eq:hlce}
\end{equation}

\notesec{Step 6, the point prediction}
The prediction scored by every correlation is the expectation over the nominal centers,
\begin{equation}
\hat z = \sum_{u=1}^{n_{\mathrm{bin}}} \hat\omega_u\, \bar\upsilon_u \in [-3.875,\ 3.875],
\label{eq:hlpoint}
\end{equation}
a convex combination of the centers, so it lies inside the outermost centers and moves
toward the middle of the histogram when the predicted mass spreads. The trainer also logs
the entropy of $\hat\omega$ and the fraction of measured entries whose argmax bin is within
one bin of the target's.

\notesec{Step 7, the pinball arm}
The quantile arm emits $n_{\mathrm{out}} = n_q = 19$ values at levels $\gamma_u = u/20$
and averages the check loss,
\begin{equation}
\mathcal{L}^{\mathrm{pin}}(\hat{\mathbf{o}}, z) = \frac{1}{n_q} \sum_{u=1}^{n_q} \max\!\big(\gamma_u (z - \hat o_u),\ (\gamma_u - 1)(z - \hat o_u)\big),
\qquad
\hat z = \hat o_{10},
\label{eq:pinball}
\end{equation}
the point being the median quantile ($\gamma_{10} = 0.5$). Nothing orders the 19 outputs,
so quantile crossing is not prevented.

\notesec{Step 8, the MSE reference}
With $n_{\mathrm{out}} = 1$, $\mathcal{L}^{\mathrm{MSE}}(\hat o, z) = (\hat o - z)^2$ and
$\hat z = \hat o$. The overnight queue dropped its run for time (Sec.~\ref{sec:measured}).

\subsection{From elements to the batch objective\secstatus{todo}}
\label{sec:batchobjective}

\notesec{Step 9, the mean within a dataset}
Write $\mathcal{L}^{\circ}$ for whichever elementwise loss the arm uses. For dataset $s$ in
batch $\mathcal{B}$, the loss is the mean over its measured entries, pooled over its records,
\begin{equation}
\mathcal{L}_s(\mathcal{B}) = \frac{\sum_{b \in \mathcal{B},\, s(b) = s}\ \sum_{i} m_{b,i}\, \mathcal{L}^{\circ}(\hat{\mathbf{o}}_{b,i}, z_{b,i})}{\sum_{b \in \mathcal{B},\, s(b) = s}\ \sum_{i} m_{b,i}},
\label{eq:dsmean}
\end{equation}
where $i$ runs over the $N$ genes of a gene-level record and over its single entry for a
record-level one. Within a dataset a record with more measured genes therefore counts for
more.

\notesec{Step 10, the weighting across datasets}
With $\mathcal{S}_{\mathcal{B}}$ the datasets that have a measured entry in the batch,
\begin{equation}
\mathcal{L}_y(\mathcal{B}) = \sum_{s \in \mathcal{S}_{\mathcal{B}}} w_s(\mathcal{B})\, \mathcal{L}_s(\mathcal{B}),
\qquad
w_s(\mathcal{B}) =
\begin{cases}
1 / |\mathcal{S}_{\mathcal{B}}| & \text{equal weighting (main)},\\
n_s(\mathcal{B}) \big/ \sum_{s' \in \mathcal{S}_{\mathcal{B}}} n_{s'}(\mathcal{B}) & \text{record weighting (arm)}.
\end{cases}
\label{eq:objective}
\end{equation}
Against the manuscript's $\mathcal{L}_y = \sum_t w_t \sum_b m^{(b)}_t \ell_t$, two things
differ: each dataset's term is normalized by its own measured count, and the weights are not
constants but functions of what the batch happens to hold. 043 adds no graph term, so
$\mathcal{L} = \mathcal{L}_y$.

\subsection{What a batch holds\secstatus{todo}}
\label{sec:batch}

\notesec{Step 11, the records sampler}
The default sampler shuffles the $n = \tdNTrain$ training records each epoch and cuts the
order into batches of $B = 8$, so a batch is a draw without replacement and its composition
$(n_s(\mathcal{B}))_{s}$ is multivariate hypergeometric (the last batch of an epoch holds 7
records, $\tdNTrain = 6{,}345 \times 8 + 7$; it is ignored below). Vaishnav holds
$n_{\mathrm{vai}} = \tdNVaishnavTrain$ of them, \tdVaishnavPct\%, so a batch holds
$8 \times (1 - n_{\mathrm{vai}}/n) = \tdGeneRecordsPerBatch$ gene-level records on average.

\notesec{Step 12, the expected weight as a gradient share}
The batch gradient is $\nabla \mathcal{L}_y = \sum_{s} w_s(\mathcal{B}) \nabla
\mathcal{L}_s(\mathcal{B})$, so the expected weight
\begin{equation}
\bar w_s = \mathbb{E}_{\mathcal{B}}\big[w_s(\mathcal{B})\big] =
\sum_{(n_1, \ldots, n_9)} \frac{\prod_{s'} \binom{n_{s'}}{n_{s'}(\mathcal{B})}}{\binom{n}{B}}\; w_s(\mathcal{B}),
\qquad
\bar w_s^{\,\mathrm{records}} = \frac{n_s}{n},
\label{eq:share}
\end{equation}
is the expected share of the objective's gradient that dataset $s$ carries, if the
per-dataset gradients $\nabla \mathcal{L}_s$ have comparable norms. The sum runs over every
composition of a batch of 8 into the nine datasets, and the record-weighting value follows
from linearity. Table~\ref{tab:share} gives both exactly. Under equal weighting Vaishnav
carries $\bar w_{\mathrm{vai}} = \tdVaishnavWeightEqual$: it is present in nearly every
batch, and the batches that also hold a gene-level record split the weight only two or three
ways. Kemmeren is present in a fraction \tdKemmerenPresent{} of batches and carries
$\bar w = \tdKemmerenWeightEqual$. Hypothesis (untested): the per-dataset gradient norms are
not comparable, since the 32-bin cross-entropy of one Vaishnav scalar and the mean over
about 6{,}000 genes of a Kemmeren profile differ in variance; the norms were not logged.

\notesec{Step 13, the balanced sampler}
The alternative draws, for each of the 8 slots of a batch, a dataset uniformly among the
nine and then a record uniformly within it, with replacement, for the same 6{,}346 batches
per epoch. The composition is then multinomial with equal cell probabilities, and both
weightings give $\bar w_s = \tdBalancedWeight$ for every dataset by symmetry. Each dataset
is drawn about $6{,}346 \times 8 / 9 \approx 5{,}641$ times per epoch, so the 33 turnover
training records are each revisited about 171 times per epoch and Vaishnav's 42{,}694 are
subsampled to about 13\%. The sampler acts on the batch and the weighting on the loss; the
two compose, and a balanced run is labeled \texttt{<arm>\_balanced\_<loss>\_<weighting>}.

%% SOURCE: experiments/043-expression-proteome-morphology/scripts/training_design_043_tables.py
%% -> tables/t2-batch-share-rows.tex (train_043_prep/prep_v1_routeA/meta.json, job 6134;
%% conf/train_043_default.yaml)
\begin{table}[htbp]
\centering
\footnotesize
\setlength{\tabcolsep}{3pt}
\caption[]{\textbf{Training records and the expected weight of each dataset in
the batch objective, \texttt{joint\_9}, batches of 8.} Records from the prepared store (job
6134). Validation records are those the trainer scores per evaluation (at most 2{,}000 per
dataset). Expected weights are exact (Eq.~\ref{eq:share}): the records sampler is a draw
without replacement (multivariate hypergeometric over all 12{,}870 compositions of 8 into
nine datasets), the balanced sampler a multinomial with probability $1/9$ per dataset.
Probabilities and weights are per batch.}
\label{tab:share}
\begin{tabular}{>{\raggedright\arraybackslash}p{26mm}>{\raggedright\arraybackslash}p{9mm}>{\raggedleft\arraybackslash}p{14mm}>{\raggedleft\arraybackslash}p{19mm}>{\raggedleft\arraybackslash}p{15mm}>{\raggedleft\arraybackslash}p{18mm}>{\raggedleft\arraybackslash}p{19mm}>{\raggedleft\arraybackslash}p{19mm}>{\raggedleft\arraybackslash}p{19mm}}
\toprule
Dataset & Level & Training records, $n_s$ & Validation records scored, n & Share of training records, \% & Probability present in a batch, records sampler & Expected weight $\bar w_s$, records sampler, equal weighting & Expected weight $\bar w_s$, records sampler, record weighting & Expected weight $\bar w_s$, balanced sampler, either weighting \\
\midrule
%% GENERATED by experiments/043-expression-proteome-morphology/scripts/training_design_043_tables.py
%% SOURCE: train_043_prep/prep_v1_routeA/meta.json (job 6134), conf/train_043_default.yaml. Do not edit by hand.
Kemmeren 2014 & gene & 1{,}202 & 140 & 2.4 & 0.174 & 0.066 & 0.024 & 0.111 \\
Messner 2023 & gene & 3{,}712 & 455 & 7.3 & 0.455 & 0.188 & 0.073 & 0.111 \\
Zelezniak 2018 & gene & 79 & 9 & 0.2 & 0.012 & 0.004 & 0.002 & 0.111 \\
Caudal 2024 & gene & 566 & 189 & 1.1 & 0.086 & 0.032 & 0.011 & 0.111 \\
Muenzner 2024 protein & gene & 476 & 159 & 0.9 & 0.073 & 0.027 & 0.009 & 0.111 \\
Muenzner 2024 turnover & gene & 33 & 11 & 0.1 & 0.005 & 0.002 & 0.001 & 0.111 \\
Boocock 2025 pooled & gene & 235 & 79 & 0.5 & 0.036 & 0.013 & 0.005 & 0.111 \\
N'Guessan 2025 & gene & 1{,}770 & 591 & 3.5 & 0.247 & 0.095 & 0.035 & 0.111 \\
Vaishnav 2022 & record & 42{,}694 & 2{,}000 of 14{,}232 & 84.1 & 1.000 & 0.573 & 0.841 & 0.111 \\
\midrule
All nine &  & 50{,}767 & 3{,}633 & 100.0 &  & 1.000 & 1.000 & 1.000 \\
\bottomrule

\end{tabular}
\end{table}

\subsection{Optimization and the validation loss\secstatus{todo}}
\label{sec:optim}

\notesec{Step 14, the optimizer}
AdamW at learning rate $3 \times 10^{-4}$ with weight decay $10^{-8}$, linear warmup over
500 steps then constant, gradient norm clipped at 10, dropout 0.1, bf16 autocast on the
card. A non-finite gradient norm stops the run.

\notesec{Step 15, the validation loss}
The validation pass computes the same objective with the same weighting, but accumulates the
per-dataset sums of Eq.~\ref{eq:dsmean} over the whole pass before combining them, so under
equal weighting
\begin{equation}
\mathcal{L}^{\mathrm{val}} = \frac{1}{|\mathcal{S}|} \sum_{s \in \mathcal{S}} \frac{\sum_{b \in \mathrm{val}_s} \sum_i m_{b,i}\, \mathcal{L}^{\circ}(\hat{\mathbf{o}}_{b,i}, z_{b,i})}{\sum_{b \in \mathrm{val}_s} \sum_i m_{b,i}},
\label{eq:valloss}
\end{equation}
logged as \texttt{val/loss} with each dataset's term as \texttt{val/<dataset>/loss}. The
training curve \texttt{train/loss} is the mean of $\mathcal{L}_y(\mathcal{B})$ over the
epoch's batches, whose weights follow Table~\ref{tab:share}; the two are the same loss but
not the same weighting of datasets, so their gap mixes generalization with the batch
composition. The validation loss was added on 2026-10-23 (commit \texttt{e6ad22d0c}); none
of the overnight runs of 2026-10-22 has it, and the prepare-experiment gate
(\file{.claude/skills/prepare-experiment/SKILL.md}) now refuses a chain whose smoke run
lacks \texttt{val/loss} and the per-dataset pair.
